% !TeX root = ../../Book3.tex
% 纲要标题：### 12.6 一阶 Cauchy–Riemann 型算子
% 纲要位置：几何学计划纲要.md，第 2592 行。
% 正文状态：未撰写；本文件目前是框架。

% 纲要内容开始
% 
% * 通过椭圆主符号获得局部 W¹ᵖ 估计
% * p>2 时的连续性及非线性复合
% * 后续闭曲面与条带的不同全局条件
% 
% 纲要内容结束

\section{一阶 Cauchy–Riemann 型算子}
\label{g3:sec:1206}

待撰写

\subsection{\texorpdfstring{通过椭圆主符号获得局部 \(W^{1,p}\) 估计}{通过椭圆主符号获得局部 W¹ᵖ 估计}}
\label{g3:subsec:1206:01}
待撰写

\subsection{\texorpdfstring{\(p>2\) 时的连续性及非线性复合}{p>2 时的连续性及非线性复合}}
\label{g3:subsec:1206:02}
待撰写

\subsection{后续闭曲面与条带的不同全局条件}
\label{g3:subsec:1206:03}
待撰写
% 纲要条目排为编号小节标题，并进入目录与书签；说明段仍保留于纲要注释中。
